% ============================================================
% 第 1 章 磁路和磁性材料（教材第一篇，印刷页 4--20）
% 节名与教材一致：第一节 磁场中基本物理量和磁路的概念；第二节 磁路的基本定律；
% 第三节 磁性材料及其特性；第四节 简单磁路计算；第五节 交流磁路中的励磁电流和磁通
% ============================================================
\chapter{磁路和磁性材料}

以电磁感应定律为基础进行能量转换或传递的机电装置主要包括：将机械能转换为电能的发电机，将电能转换为机械能的电动机，以及实现电能传递的变压器。这些装置在能量传递或转换过程中均以\textbf{电磁场}为媒介，因此电磁场的强弱和分布不仅影响电机与变压器的性能，也决定着它们的体积和质量。

电磁场问题的求解一般有两种方法：一是\textbf{场的分析方法}，在给定区域用离散点物理量的数值定义其电磁特性，计算准确但计算难度与计算量较大，结构复杂时往往难以完成；二是\textbf{路的分析方法}，从宏观角度出发，把给定区域的场量与电路中的物理量作类比，用集总参数等效离散参数，把磁场问题简化成磁路问题。路的方法准确度低于场的方法，但计算简单，工程上足够用，因此本书主要采用路的分析方法。

\section{磁场中基本物理量和磁路的概念}

\subsection{磁场中常用物理量}

\subsubsection{磁感应强度}
当闭合导体中通入电流时，导体周围就会产生磁场，磁场的强弱和方向用\textbf{磁感应强度} $\boldsymbol{B}$ 来描述，其大小可由毕奥—萨伐尔定律求得。设导体中电流为 $i$，每一线元 $\mathrm{d}\boldsymbol{l}$ 在距其 $R$ 处的点 $P$ 产生的磁感应强度元为
\begin{equation}
\mathrm{d}\boldsymbol{B}=\frac{\mu_0}{4\pi}\cdot\frac{i\,\mathrm{d}\boldsymbol{l}\times \boldsymbol{a}_R}{R^{2}},
\end{equation}
式中 $\boldsymbol{a}_R$ 为由电流元指向点 $P$ 的单位矢量，$R$ 为从电流元到点 $P$ 的距离，$\mu_0$ 为真空磁导率。

$\boldsymbol{B}$ 的方向也可用闭合的磁感应矢量线（磁力线）描述：通电导体中的电流与所产生的磁场磁力线之间符合\textbf{右手螺旋定则}。$\boldsymbol{B}$ 是矢量，国际单位制中单位为 T（特斯拉）。

\subsubsection{磁通量}
$\boldsymbol{B}$ 描述的是空间每一点上磁场的大小和方向；要描述一个给定面上的磁场情况，需引入\textbf{磁通量}（简称磁通）$\Phi$。当磁感应强度在给定开表面 $A$ 上不均匀分布时
\begin{equation}
\Phi=\int_{A}\boldsymbol{B}\cdot \mathrm{d}\boldsymbol{A}.
\end{equation}
当 $\boldsymbol{B}$ 均匀分布且磁力线垂直穿过表面时
\begin{equation}
\Phi=BA.
\end{equation}

$\Phi$ 的单位为 Wb（韦伯），$A$ 的单位为 $\mathrm{m^2}$，故 $1\,\mathrm{T}=1\,\mathrm{Wb/m^2}$。工程上有时用 Gs（高斯）与 Mx（麦克斯韦）：
\begin{equation}
1\,\mathrm{Mx}=10^{-8}\,\mathrm{Wb},\qquad 1\,\mathrm{Gs}=10^{-4}\,\mathrm{T}.
\end{equation}
由 $\Phi=BA$ 可见，磁感应强度可看成穿过单位面积的磁通量，故又称\textbf{磁通密度}（磁密）。

\subsubsection{磁导率}
通电导体所产生磁场的强弱与导体周围介质的磁性能密切相关：有些介质使磁场显著增强，另一些则使其略有削弱。表征介质磁性能的物理量称为\textbf{磁导率} $\mu$，单位 H/m。真空磁导率
\begin{equation}
\mu_0=4\pi\times10^{-7}\ \mathrm{H/m},
\end{equation}
铁磁材料的磁导率 $\mu\gg\mu_0$。磁场计算时常采用\textbf{相对磁导率} $\mu_r=\mu/\mu_0$：各种硅钢片材料 $\mu_r$ 可达 5000、6000；铸钢约为 1000。

\subsubsection{磁场强度}
表征磁场性能的另一个基本物理量是\textbf{磁场强度} $\boldsymbol{H}$，定义为
\begin{equation}
\boldsymbol{H}=\frac{\boldsymbol{B}}{\mu}\qquad\text{或}\qquad \boldsymbol{B}=\mu\boldsymbol{H}.
\end{equation}
磁场强度反映了单位长度磁路上所消耗的磁动势，又称单位长度上的磁压降。$\boldsymbol{H}$ 也是矢量，单位为 A/m。

\subsection{磁路的概念}
电路中把电流流过的路径称为电路；与之相类比，研究磁场问题时把\textbf{磁通所经过的路径}称为\textbf{磁路}。电机和变压器中，为使磁通按设定路径通过，磁路一般由铁磁材料组成。

磁路可类比于电路，但两者有本质区别：
\begin{itemize}
  \item 电路中的电流由带电粒子的定向运动产生；磁路中\textbf{没有任何物质沿着闭合回路流动}。
  \item 物质世界存在电的导体与绝缘体，电流只在导体中流动；但\textbf{不存在磁的绝缘体}，磁通可以存在于任何介质之中。
\end{itemize}

由于铁磁材料磁导率远大于空气，绝大部分磁通在铁芯内通过——铁芯路径是设定的磁路，其中的磁通称为\textbf{主磁通}，其路径称为\textbf{主磁路}；铁芯外非铁磁材料磁导率不为零，总有少量磁通经铁芯外路径闭合，称为\textbf{漏磁通}，其路径称为\textbf{漏磁路}。电机和变压器中能量的转换与传递通过主磁通完成，故主磁通又称\textbf{工作磁通}。

产生磁通的载流线圈称为\textbf{励磁线圈}（励磁绕组），其中电流称为\textbf{励磁电流}。直流励磁电流产生的磁通不随时间变化，相应的磁路称为\textbf{直流磁路}；交流励磁电流产生的磁通随时间变化，称为\textbf{交流磁路}。

\begin{center}
\begin{tabular}{@{}p{2.4cm}p{4.9cm}p{5.6cm}@{}}
\toprule
对比项 & 电路 & 磁路 \\
\midrule
“流动”的物理量 & 电流 $I$（电荷定向运动） & 磁通 $\Phi$（无物质流动） \\
“源” & 电动势 $E$ & 磁动势 $F=Ni$ \\
“阻力” & 电阻 $R=l/(\gamma A)$ & 磁阻 $R_m=l/(\mu A)$ \\
介质界限 & 有导体与绝缘体之分 & 无磁绝缘体，任何介质均可通磁 \\
漏泄 & 漏电通常可忽略 & 漏磁严重，精确计算必须考虑 \\
线性 & 金属电阻基本为常数 & $\mu$ 随饱和程度变化，$R_m$ 非线性 \\
叠加原理 & 线性电路可用 & 非线性磁路\textbf{不可}用 \\
\bottomrule
\end{tabular}
\end{center}

\corepoint{磁路与电路只是"形式类比"，不是"物理同构"。最容易错的两处：① 磁路中 $\mu$ 不是常数，因此不能像电路那样对多个磁动势用叠加原理；② 漏磁不能忽略，尤其在气隙附近。做题时若发现要用叠加，先检查磁路是否处于线性段（$\mu_r$ 恒定）——只有线性化之后才可以。}

\section{磁路的基本定律}

\subsection{安培环路定律}
沿着任何一条闭合回路 $l$，磁场强度 $\boldsymbol{H}$ 的线积分等于该闭合回路所包围的总电流（代数和）：
\begin{equation}
\oint_{l}\boldsymbol{H}\cdot \mathrm{d}\boldsymbol{l}=\sum i .
\end{equation}
其中，若电流正方向与回路上磁场强度的方向符合右手螺旋关系，电流取正值，否则取负值。若沿回路 $l$ 磁场强度大小处处相等，且闭合回路所包围的总电流由通有电流 $i$ 的 $N$ 匝线圈提供，当 $Hl$ 的回路方向与电流方向符合右手螺旋定则时，可简写为
\begin{equation}
Hl=Ni .
\end{equation}

\subsection{磁路的欧姆定律}
设一等截面无分支铁芯磁路，铁芯截面积 $A$、磁路平均长度 $l$、材料磁导率 $\mu$，励磁绕组匝数 $N$、电流 $i$。若不计漏磁通，并认为各截面上磁密均匀分布且垂直于截面，则
\begin{equation}
\Phi=BA=\mu H A=\mu\frac{Ni}{l}A=\frac{Ni}{l/(\mu A)}=\frac{F}{R_m},
\end{equation}
其中
\begin{equation}
F=Ni,\qquad R_m=\frac{l}{\mu A},\qquad \Lambda=\frac{1}{R_m}.
\end{equation}

\keyformula{\Phi=\frac{F}{R_m}=\frac{Ni}{l/(\mu A)}}

式中：$F$ 为作用在铁芯磁路上的安匝数，称为\textbf{磁动势}，单位 A；$R_m$ 为\textbf{磁阻}，单位 A/Wb 或 $1/\mathrm{H}$；$\Lambda$ 为磁阻的倒数，称为\textbf{磁导}，单位 Wb/A 或 H。该式与电路欧姆定律形式相似，故称\textbf{磁路的欧姆定律}：$F$ 类比于 $E$，$\Phi$ 类比于 $I$，$R_m$ 类比于 $R$。

磁阻主要取决于磁路的几何尺寸和材料磁导率。由于铁磁材料 $\mu$ 不是常数，\textbf{铁磁材料的磁阻是非线性的}，这正是磁路与电路的本质区别之一。

\subsection{磁路的基尔霍夫第一定律}
若铁芯是带分支的磁路，不计漏磁时，任取一闭合面 $A$，由磁通连续性原理，穿过闭合面磁通的代数和必为零——进入该闭合面的磁通等于离开该闭合面的磁通：
\begin{equation}
\Phi_1=\Phi_2+\Phi_3
\quad\Longleftrightarrow\quad
-\Phi_1+\Phi_2+\Phi_3=0
\quad\Longleftrightarrow\quad
\sum\Phi=0 .
\end{equation}
规定穿出闭合面的磁通为正、进入为负，即得磁路的基尔霍夫第一定律。

\subsection{磁路的基尔霍夫第二定律}
磁路计算中通常把磁路分段：凡是材料和横截面积均相同、且磁通也相等的磁路作为一段。每一段内各点磁场强度相同。对由 $k$ 段组成的闭合磁路，由安培环路定律
\begin{equation}
H_1l_1+H_2l_2+H_3l_3=Ni
\quad\Longleftrightarrow\quad
\sum_{k}H_k l_k=Ni .
\end{equation}
式中，若 $H_k$ 的正方向与闭合回路方向一致，$H_kl_k$ 前取“$+$”号，否则取“$-$”号；若电流 $i$ 的正方向与 $l$ 的方向符合右手螺旋关系，$Ni$ 前取“$+$”号，否则取“$-$”号。

\keyformula{\sum_{k}H_k l_k=Ni}

$H_kl_k$ 称为第 $k$ 段磁路上的\textbf{磁位降}，$Ni$ 为作用在磁路上的\textbf{总磁动势}。该式表明：作用在任何闭合磁路的总磁动势恒等于各段磁路磁压降的代数和——即磁路的基尔霍夫第二定律，它实质上是安培环路定律的另一种表达形式。

\begin{center}
\begin{tabular}{@{}lll@{}}
\toprule
磁路定律 & 表达式 & 电路中的对应式 \\
\midrule
安培环路定律 & $\oint \boldsymbol{H}\cdot\mathrm{d}\boldsymbol{l}=\sum i$ & 全电流定律 \\
磁路欧姆定律 & $\Phi=F/R_m$，$R_m=l/(\mu A)$ & $I=E/R$ \\
基尔霍夫第一定律 & $\sum\Phi=0$ & $\sum I=0$ \\
基尔霍夫第二定律 & $\sum H_k l_k=Ni$ & $\sum U=\sum E$ \\
\bottomrule
\end{tabular}
\end{center}

\corepoint{用第二定律时符号最易错。稳妥做法：先任意规定回路绕行正方向，再逐段判断 $H_k$ 与绕行方向是否一致、$Ni$ 与绕行方向是否成右螺旋，同号取正、异号取负，最后解代数方程。切忌"全取正号"硬算。}

\section{磁性材料及其特性}

磁性材料按磁性能分为\textbf{顺磁材料}、\textbf{反磁材料}和\textbf{铁磁材料}。顺磁材料 $\mu$ 略大于 $\mu_0$，反磁材料 $\mu$ 略小于 $\mu_0$；工程上二者的磁导率一般按等于 $\mu_0$ 计算，如空气、变压器油、铜和铝等。铁磁材料 $\mu\gg\mu_0$，一般大到几千倍，如铁、镍、钴及其合金。电机和变压器的磁路通常都由铁磁材料制成。

\subsection{铁磁材料的磁化}
铁磁材料放入磁场中会呈现很强的磁性，这种现象称为\textbf{磁化}，能被磁化是铁磁材料的基本特性之一。普遍认为，铁磁材料内部存在许多很小的、被称为\textbf{磁畴}的天然磁化区，每个磁畴可看做一个微型磁针。无外磁场时各磁畴任意排列，材料不显磁性；放入外磁场后，磁畴受磁力作用发生旋转并沿磁场方向排列整齐，材料显示较强磁性。由此形成的磁化磁场与外磁场叠加，使合成磁场显著增强。非铁磁材料中不含磁畴，故在同样磁场强度下磁场不会显著增强——这就是铁磁材料磁导率大得多的原因。

\subsection{铁磁材料的磁化曲线}

\subsubsection{原始磁化曲线}
将未磁化的铁磁材料磁化，$H$ 由零逐渐增加时 $B$ 随之增加，所得 $B=f(H)$ 曲线称为\textbf{原始磁化曲线}。其形状分四段：起始段（$B$ 随 $H$ 增加较慢）、直线段（$B$ 迅速增长）、饱和段（增长放慢并逐渐趋于饱和）、过饱和段（曲线基本成为与非铁磁材料 $B=\mu_0H$ 平行的直线）。$b$ 点称\textbf{膝点}，$c$ 点称\textbf{饱和点}。

由于磁化曲线不是直线，$\mu=B/H$ 也随 $H$ 变化：直线段内 $\mu$ 随 $H$ 增加迅速增大；进入饱和区后 $\mu$ 急剧下降；至过饱和区 $\mu$ 基本不变，保持为常数。因此\textbf{铁磁材料的磁阻不是常数，而随饱和程度增加而增大}。

在电机和变压器的磁路设计中，为了用较少的磁动势获得较大的磁通，通常把铁磁材料的工作点选在\textbf{膝点附近}。

\subsubsection{磁滞回线}
将铁磁材料在 $-H_m$ 到 $+H_m$ 之间反复周期性磁化，得到对称于原点的闭合曲线，即\textbf{磁滞回线}。当 $H$ 从零增到 $+H_m$，$B$ 增到 $B_m$；再逐渐减小 $H$，$B$ 沿曲线下降，当 $H=0$ 时 $B$ 并不为零，而是等于 $B_r$——去掉外加磁场后材料内部仍保留一定磁感应强度的现象称为\textbf{顽磁性}，$B_r$ 称为\textbf{剩余磁感应强度}（剩磁）。要使 $B$ 从 $B_r$ 减小到零，必须加反向磁场，该反向磁场强度称为\textbf{矫顽力} $H_c$。$B_r$ 与 $H_c$ 是铁磁材料的两个重要参数。

$B$ 的变化总是落后于 $H$ 的变化，这种现象称为\textbf{磁滞}，相应的 $B$—$H$ 闭合曲线称为磁滞回线。磁滞现象是铁磁材料的另一个基本特性。

\subsubsection{基本磁化曲线}
对同一材料用不同的 $H_m$ 周期性磁化，可得一系列大小不同的磁滞回线；将各磁滞回线的\textbf{顶点}连接起来所得的曲线称为\textbf{基本磁化曲线}。工程上进行磁路计算时采用的都是基本磁化曲线。

\subsection{铁磁材料的分类}
按磁滞回线形状，铁磁材料分为两类：
\begin{itemize}
  \item \textbf{软磁材料}：磁滞回线窄、剩磁和矫顽力都小，如铸铁、铸钢、硅钢片和坡莫合金等。磁滞损耗小、磁导率高，常用来制造变压器和电机的铁芯。
  \item \textbf{硬磁材料}：磁滞回线宽、剩磁和矫顽力都大，如铝镍钴合金和稀土合金等。因剩磁大，常用来制造永久磁铁，故又称\textbf{永磁材料}。
\end{itemize}

\subsection{铁芯损耗}
\textbf{铁芯损耗}（铁耗）是铁芯磁路中由于磁滞和涡流引起的功率损耗的总称，\textbf{只存在于交流磁路中}，直流磁路没有这项功率损耗。

\subsubsection{磁滞损耗}
外磁场作用下，磁畴为保持磁场方向与外磁场一致会发生转动；外磁场交变时磁畴来回翻转、彼此摩擦而引起损耗，称为\textbf{磁滞损耗}。实验证明，磁滞损耗正比于交变磁场频率 $f$、磁滞回线面积和铁磁材料体积 $V$，而磁滞回线面积又正比于最大磁感应强度 $B_m$ 的 $n$ 次方（$n$ 取决于材料，一般硅钢片 $n=1.6\sim2.3$）：
\begin{equation}
P_h=C_h f B_m^{n} V .
\end{equation}

\subsubsection{涡流损耗}
磁通交变时，由电磁感应定律，铁芯内将产生感应电动势和电流，这些电流在铁芯内部围绕磁通形成环流，称为\textbf{涡流}；涡流在铁芯中引起的电阻损耗称为\textbf{涡流损耗}：
\begin{equation}
P_e=C_e f^{2}B_m^{2}\delta^{2} V ,
\end{equation}
式中 $C_e$ 为涡流损耗系数，反比于材料电阻率；$\delta$ 为硅钢片厚度。该式表明：为减小涡流损耗，可在钢材中加入少量硅以增加电阻率，并用相互绝缘的许多薄硅钢片叠成铁芯代替整块铁芯——变压器和电机的铁芯就是采用硅钢片制作的。

\subsubsection{铁芯损耗}
将磁滞损耗与涡流损耗相加即得铁芯损耗：
\begin{equation}
P_{Fe}=P_h+P_e=\left(C_h f B_m^{n}+C_e f^{2}B_m^{2}\delta^{2}\right)V .
\end{equation}
对于一般硅钢片，在正常工作磁感应强度范围内（$1\,\mathrm{T}<B_m<1.8\,\mathrm{T}$），可近似写成
\begin{equation}
P_{Fe}\approx C_{Fe}f^{1.3}B_m^{2}G ,
\end{equation}
式中 $C_{Fe}$ 为铁芯损耗系数，$G$ 为铁芯质量。

\keyformula{P_{Fe}\approx C_{Fe}f^{1.3}B_m^{2}G}

该式表明：铁芯损耗近似与频率的 1.3 次方成正比，与铁芯中最大磁感应强度的二次方及铁芯质量成正比。

\begin{center}
\begin{tabular}{@{}llll@{}}
\toprule
损耗 & 成因 & 表达式 & 减小措施 \\
\midrule
磁滞损耗 & 磁畴反复翻转、相互摩擦 & $P_h=C_hfB_m^{n}V$ & 选软磁材料（硅钢片、坡莫合金） \\
涡流损耗 & 交变磁通在铁芯内感生环流 & $P_e=C_ef^{2}B_m^{2}\delta^{2}V$ & 加硅提高电阻率、薄片叠装、片间绝缘 \\
铁芯损耗 & 以上两者之和 & $P_{Fe}\approx C_{Fe}f^{1.3}B_m^{2}G$ & 降低 $B_m$、减薄片厚、提高电阻率 \\
\bottomrule
\end{tabular}
\end{center}

\corepoint{铁耗 $\propto B_m^2$、$\propto f^{1.3}$，这条比例关系贯穿全书：变压器空载损耗近似等于铁耗（与负载无关），而铜耗随负载电流平方变化；感应电机里铁耗基本为常数、被归入空载损耗。凡遇到"某量随电压/频率如何变化"的判断，先想 $B_m\propto U/f$ 这一层关系。}

\section{简单磁路计算}

磁路计算问题一般分为两类。

\begin{enumerate}
  \item \textbf{给定磁通求磁动势}（正问题）。磁路材料和尺寸给定，要求产生给定磁通所需的励磁磁动势。由于铁磁材料非线性，必须先用 $\Phi$ 与各段有效面积求出各段 $B$，再对每个 $B$ 在基本磁化曲线上查出相应 $H$，最后由 $\sum H_kl_k=Ni$ 求出 $F$ 与励磁电流。
  \item \textbf{给定磁动势求磁通}（反问题）。由于 $\mu\neq$ 常数，通常先\textbf{假设}磁通，按正问题算出所需磁动势，与给定值比较后修正假设，用\textbf{迭代法}逐步逼近。
\end{enumerate}

\subsection{计算要点与步骤}
\begin{enumerate}
  \item \textbf{分段}。按"材料相同、截面积相同且磁通相等"把磁路分成若干段；有气隙时气隙单独成段。
  \item \textbf{求各段中心线长度 $l_k$ 与有效截面积 $A_k$}。硅钢片叠装时须乘\textbf{叠片系数}（除去片间绝缘厚度），即 $A=A_{\text{几何}}\times K_{Fe}$。
  \item \textbf{气隙有效面积}。考虑气隙磁场的边缘效应，计算气隙有效面积时通常在长、宽方向各增加一个气隙长度，即 $A_\delta=(a+\delta)(b+\delta)$。
  \item \textbf{求各段磁密} $B_k=\Phi/A_k$。
  \item \textbf{求各段磁场强度} $H_k$：线性段可用 $H_k=B_k/(\mu_r\mu_0)$；非线性段须查基本磁化曲线。
  \item \textbf{求磁压降与磁动势}：$F_k=H_kl_k$（或 $F_k=R_{mk}\Phi$），$F=\sum F_k$，励磁电流 $I=F/N$。
\end{enumerate}

\subsection{两种算式的选用}
\begin{itemize}
  \item \textbf{$F=Hl$ 法}：适用于材料磁导率非线性、需查磁化曲线的场合，是工程计算的主方法。
  \item \textbf{$F=R_m\Phi$ 法}：只在 $\mu$ 为常数的线性磁路（如气隙、或明确给出 $\mu_r$ 恒定的铁芯）中方便。
\end{itemize}

一个必须记住的结论：\textbf{气隙虽短，磁压降却占绝大部分}。教材算例中气隙仅占磁路平均长度的 0.36\%，其磁阻与磁压降却占整个磁路的 92\%；即使铁芯部分已进入饱和、该段磁阻增大，气隙磁压降仍占励磁磁动势的一半以上。

\corepoint{三处高频错误：① 忘记乘叠片系数，导致 $B$ 偏小、$H$ 查表偏小、$F$ 偏小；② 气隙面积直接取几何面积，忘了长宽各加一个 $\delta$，使气隙磁阻偏大；③ 铁芯进入饱和后仍用常数 $\mu_r$ 计算，结果严重偏小（饱和段 $\mu$ 急剧下降，$H$ 会跳升一到两个数量级）。}

\section{交流磁路中的励磁电流和磁通}

交流铁芯磁路中，由于铁芯磁化曲线的非线性，励磁电流 $i_m$ 与主磁通的波形会产生\textbf{畸变}。以单一铁芯材料的交流磁路为例，为便于图解，把磁化特性写成 $\Phi=f(i)$ 的形式。

\subsection{磁化特性为基本磁化曲线}
磁路内磁通较小时，磁路不饱和，$\Phi$ 与 $i$ 基本为线性关系：若磁通随时间按正弦规律变化，励磁电流也按正弦规律变化，两者\textbf{同相位}，波形相同。

磁通较大时磁路饱和，$\Phi$ 与 $i$ 呈非线性，需用图解法分析，分两种情况。

\subsubsection{磁通为正弦波时励磁电流的波形}
设磁通随时间作正弦变化，取若干时刻的磁通瞬时值，由 $\Phi=f(i)$ 曲线查出产生该磁通所需的励磁电流，逐点描出即得 $i_m=f(t)$。可见：磁通正弦且磁路饱和时，由于非线性，\textbf{励磁电流波形发生畸变，成为尖顶波}。分解后除基波外还含其他奇次谐波，其中\textbf{3 次谐波最大}。磁路越饱和，尖顶越严重，谐波越显著。但无论尖顶多严重，\textbf{其基波相位始终与磁通同相}。

\subsubsection{励磁电流为正弦波时磁通的波形}
当 $i_m$ 随时间正弦变化且磁路饱和时，磁通波形也发生畸变，成为\textbf{平顶波}；分解后除基波外亦含奇次谐波，同样以 3 次谐波最大。磁路越饱和，平顶越严重。但无论平顶多严重，\textbf{其基波相位始终与励磁电流同相}。

综上：励磁电流与磁通波形是否畸变，决定于磁路是否饱和。不饱和时磁化曲线基本线性，两者波形非常接近；饱和后，磁通为正弦时电流畸变成尖顶波，电流为正弦时磁通畸变成平顶波。无论尖顶波还是平顶波，分解后的主要成分都是\textbf{基波和 3 次谐波}；磁通基波与励磁电流基波的相位始终相同，与是否饱和及饱和程度无关。

\subsection{磁化特性为磁滞回线}
考虑磁滞现象时，磁化特性为磁滞回线。当磁通为正弦波时，通过磁滞回线用图解法求得 $i_m=f(t)$，它为非正弦波，可分解为两个分量：
\begin{equation}
i_m=i_\mu+i_{Fe},\qquad i_{Fe}=i_h+i_e ,
\end{equation}
其中 $i_\mu$ 为与磁通同相位的尖顶波，与忽略磁滞、只考虑饱和时得到的尖顶波完全一致，即不考虑磁滞时的励磁电流，与基本磁化曲线相对应，称为\textbf{磁化电流}；$i_{Fe}$ 波形近似正弦，相位\textbf{超前磁通 $90^\circ$}，是反映磁滞损耗和涡流损耗的\textbf{有功电流}。其中 $i_h$ 是考虑了磁滞现象才存在的\textbf{磁滞损耗电流}（不计磁滞时为零），$i_e$ 是\textbf{涡流损耗电流}，两者合称\textbf{铁耗电流}。

\subsection{三点结论}
\begin{enumerate}
  \item 激励产生磁通的电流是\textbf{磁化电流} $i_\mu$，其基波与磁通同相位；由于磁路非线性，$i_\mu$ 与磁通波形不同：磁通为正弦时 $i_\mu$ 为尖顶波，电流为正弦时磁通为平顶波。
  \item 铁芯中存在磁滞与涡流损耗，励磁电流除磁化电流外还存在铁芯损耗电流，后者接近正弦，相位\textbf{超前磁通 $90^\circ$}。
  \item 励磁电流是磁化电流与铁芯损耗电流的合成，即 $i_m=i_{Fe}+i_\mu$。变压器与电机铁芯损耗很小，$i_{Fe}$ 通常仅为 $i_m$ 的 10\% 左右，故 $i_m$ 主要由 $i_\mu$ 决定。
\end{enumerate}

\corepoint{三条常考结论要能脱口而出：① 磁通正弦 $\Rightarrow$ 励磁电流尖顶波；② 励磁电流正弦 $\Rightarrow$ 磁通平顶波；③ 无论哪种情况，\textbf{磁通基波与励磁电流基波同相位}，而铁耗电流超前磁通 $90^\circ$。第 ④ 章变压器的相量图里，$I_0$ 与 $\dot\Phi_m$ 的夹角 $\alpha$（铁耗角）就是从这里来的。}

\begin{center}
\begin{tabular}{@{}llll@{}}
\toprule
磁路状态 & 给定波形 & 另一量波形 & 主要谐波 \\
\midrule
不饱和 & 磁通正弦 & 电流正弦 & 无（近正弦） \\
饱和 & 磁通正弦 & 电流尖顶波 & 3 次最大 \\
饱和 & 电流正弦 & 磁通平顶波 & 3 次最大 \\
计及磁滞 & 磁通正弦 & 电流非正弦（含 $i_\mu+i_{Fe}$） & 3 次最大；$i_{Fe}$ 超前 $90^\circ$ \\
\bottomrule
\end{tabular}
\end{center}
